\documentclass[11pt]{article}
\usepackage[a4paper,margin=1in]{geometry}
\usepackage{amsmath,amssymb,amsthm,mathtools}
\usepackage{graphicx}
\usepackage[hidelinks]{hyperref}
\usepackage{microtype}
\usepackage{enumitem}

\newtheorem{theorem}{Theorem}
\newtheorem{proposition}[theorem]{Proposition}
\newtheorem{lemma}[theorem]{Lemma}
\newtheorem{corollary}[theorem]{Corollary}
\theoremstyle{definition}
\newtheorem{definition}[theorem]{Definition}
\theoremstyle{remark}
\newtheorem{remark}[theorem]{Remark}

\newcommand{\T}{\mathbb T}
\newcommand{\Z}{\mathbb Z}
\newcommand{\R}{\mathcal R}
\newcommand{\M}{\mathcal M}
\newcommand{\A}{\mathcal A}
\newcommand{\1}{\mathbf 1}
\newcommand{\ip}[2]{\left\langle #1,#2\right\rangle}

\title{A noncommutative $2$-adic Fourier system\\
and an obstruction to the strict doubled-character model}
\author{Candidate solution packet for arXiv:2506.14431, Question 5.17}
\date{11 August 2026}

\begin{document}
\maketitle

\begin{center}
\fbox{\parbox{0.9\textwidth}{\textbf{Status.}
Candidate full dichotomy, likely valid.  The affirmative construction answers
the source question under four explicit Fourier-system axioms distilled from
the surrounding paper.  The source does not itself define the requested term,
so the semantic match remains a point for human review.}}
\end{center}

\section{Source question}

Jiao, Luo, Zhao, and Zhou ask the following on page 31, Question 5.17, of
\cite{JLZZ}:

\begin{quote}
Recall that the construction of noncommutative Walsh and Vilenkin systems
appear in the series of the papers by Sukochev and his collaborators.  It is
natural to ask whether we can construct a noncommutative $2$-adic integer
system (the classical version is in Example 2.7)?
\end{quote}

\begin{figure}[ht]
  \centering
  \includegraphics[width=\textwidth]{figures/open_problem_crop.png}
  \caption{Full-width crop of Question 5.17 on page 31 of \cite{JLZZ}.}
\end{figure}

The paper does not formally define ``noncommutative $2$-adic integer
system.''  Its noncommutative Vilenkin system has four central structural
features: it lives in the hyperfinite $\mathrm{II}_1$ factor, consists of a
complete orthonormal family of unitaries, is closed projectively under
multiplication, and is built from the classical character data.  We formulate
the answer using precisely those features.  We then prove that the stricter
version obtained by doubling the classical character group, exactly as in the
Vilenkin construction, is impossible.

\section{Definitions}

Let
\[
 K=\Z_2,\qquad A=\widehat K=\Z[1/2]/\Z.
\]
Thus $A$ is the Pr\"ufer $2$-group.  If
$r=a/2^N+\Z\in A$ and $x\in K$, define
\[
 \chi_r(x)=\exp\!\left(2\pi i\frac{a(x\bmod 2^N)}{2^N}\right).
\]
This is independent of the representative and identifies $A$ with the
character group of $K$.

\begin{definition}[Structural $2$-adic Fourier system]
A structural noncommutative $2$-adic Fourier system is a finite factor
$(\M,\tau)$ together with a countable family of unitaries
$\{X_{r,n}:r\in A,n\in\Z\}$ such that
\begin{enumerate}[label=(\roman*)]
\item the family is an orthonormal basis of $L_2(\M,\tau)$;
\item products of basis elements are scalar multiples of basis elements;
\item $\{X_{r,0}:r\in A\}$ is the classical character system of $\Z_2$;
\item $\M$ is the hyperfinite $\mathrm{II}_1$ factor.
\end{enumerate}
\end{definition}

The integer index records powers of one noncommuting translation.  The strict
doubled-character version would instead ask for a projective unitary basis
indexed by $A\times A$, mirroring the doubled dual group in the source's
noncommutative Vilenkin construction.

\section{Proof intuition}

Translation by $1$ is the odometer on $\Z_2$.  Its action on a character is
diagonal: translating $\chi_s$ multiplies it by the scalar
$e^{-2\pi i s}$.  The crossed product therefore turns the ordinary characters
into explicit noncommuting Fourier monomials $\chi_r u^n$.  Orthogonality is
read directly from the canonical crossed-product trace.

The only point not immediate from the multiplication formula is that the
ambient algebra really is the hyperfinite $\mathrm{II}_1$ factor.  Dyadic
residue classes form exact odometer towers.  Matrix units obtained by moving
between levels of a tower give nested copies of $M_{2^N}$.  Their cyclic shift
differs from the true odometer unitary only on the top level, whose trace is
$2^{-N}$; hence the difference tends to zero in $L_2$.

The strict obstruction is dual to this mechanism.  Two copies of the Pr\"ufer
group are too divisible to support a nontrivial scalar commutator pairing:
their dual is $\Z_2^2$, which contains no nonzero divisible subgroup.  The
integer translation direction in the positive construction is exactly what
escapes that obstruction.

\section{Affirmative construction}

\begin{theorem}[Odometer $2$-adic system]\label{thm:positive}
Let $\alpha(x)=x+1$ on $K=\Z_2$, and form the group-measure-space von Neumann
algebra
\[
 \M=L_\infty(K)\rtimes_\alpha\Z.
\]
Let $u$ be the canonical implementing unitary, with
$ufu^*=f\circ\alpha^{-1}$.  For $(r,n)\in A\times\Z$, set
\[
 X_{r,n}=\chi_r u^n.
\]
Then $\M$ is the hyperfinite $\mathrm{II}_1$ factor and
$\{X_{r,n}\}_{(r,n)\in A\times\Z}$ is an orthonormal basis of $L_2(\M)$.
Moreover,
\begin{equation}\label{eq:product}
 X_{r,n}X_{s,m}=e^{-2\pi i ns}X_{r+s,n+m}.
\end{equation}
The subfamily $\{X_{r,0}:r\in A\}$ is precisely the classical character
system of the $2$-adic integers, and the whole system is noncommutative.
\end{theorem}

\begin{proof}
The character identity $\chi_s(x-1)=e^{-2\pi i s}\chi_s(x)$ gives
\[
 u^n\chi_su^{-n}=e^{-2\pi i ns}\chi_s.
\]
Equation \eqref{eq:product} follows.  Taking $n=1$ and $s=1/2+\Z$ gives
$u\chi_s=-\chi_su$, so the system is genuinely noncommutative.

The canonical trace on the algebraic crossed product is
\[
 \tau\!\left(\sum_{k\in\Z}f_k u^k\right)=\int_K f_0\,d\mu.
\]
Consequently,
\[
 \ip{X_{r,n}}{X_{s,m}}
 =\tau(X_{r,n}^*X_{s,m})
 =\delta_{n,m}\int_K\overline{\chi_r}\chi_s\,d\mu
 =\delta_{n,m}\delta_{r,s}.
\]
The algebraic crossed product is $L_2$-dense in $\M$, while the characters
$\{\chi_r:r\in A\}$ are an orthonormal basis of $L_2(K)$.  Approximating each
coefficient $f_k$ by characters proves completeness of the displayed family.

We next verify the factor and hyperfinite assertions.  Translation by $1$ is
free: if $x+n=x$ in $\Z_2$ for an integer $n$, then $n=0$.  It is ergodic as
well.  Indeed, if $f\circ\alpha=f$, then its Fourier coefficient at $r$ obeys
\[
 \widehat f(r)=e^{2\pi i r}\widehat f(r),
\]
and $e^{2\pi i r}\ne1$ for every nonzero $r\in A$.  Hence $f$ is constant.
The usual Fourier-coefficient computation in the crossed product is also
short here: if a central element has coefficients $f_n$, commutation with
$L_\infty(K)$ and freeness force $f_n=0$ for $n\ne0$; commutation with $u$
then makes $f_0$ $\alpha$-invariant.  Thus the center is scalar and $\M$ is a
$\mathrm{II}_1$ factor.

For a direct proof of hyperfiniteness, fix $N\ge1$, put $d=2^N$, and let
\[
 C_j^{(N)}=j+d\Z_2,\qquad p_j^{(N)}=\1_{C_j^{(N)}}quad(0\le j<d).
\]
Define
\[
 e_{ij}^{(N)}=p_i^{(N)}u^{i-j}p_j^{(N)}.
\]
Since $u^k p_j^{(N)}u^{-k}=p_{j+k}^{(N)}$ with indices modulo $d$, these
elements obey
\[
 e_{ij}^{(N)}e_{k\ell}^{(N)}=\delta_{j,k}e_{i\ell}^{(N)},
 \qquad (e_{ij}^{(N)})^*=e_{ji}^{(N)}.
\]
They therefore span a matrix algebra $F_N\cong M_d$.  Splitting each residue
class modulo $d$ into its two residue classes modulo $2d$ gives
\begin{equation}\label{eq:nesting}
 e_{ij}^{(N)}=e_{ij}^{(N+1)}+e_{i+d,j+d}^{(N+1)},
\end{equation}
so $F_N\subset F_{N+1}$.

The diagonal of $\bigcup_NF_N$ contains every dyadic cylinder function and
is therefore $L_2$-dense in $L_\infty(K)$.  Finally define the cyclic shift
\[
 v_N=\sum_{j=0}^{d-2}e_{j+1,j}^{(N)}+e_{0,d-1}^{(N)}\in F_N.
\]
The unitaries $u$ and $v_N$ agree on $1-p_{d-1}^{(N)}$.  Hence
\[
 \|u-v_N\|_2
 =\|(u-v_N)p_{d-1}^{(N)}\|_2
 \le 2\tau(p_{d-1}^{(N)})^{1/2}
 =\frac{2}{\sqrt d}\longrightarrow0.
\]
Thus the $L_2$-closure of $\bigcup_NF_N$ contains $L_\infty(K)$ and $u$, so it
is all of $L_2(\M)$.  Therefore $\M$ is hyperfinite.  Being a separable
hyperfinite $\mathrm{II}_1$ factor, it is $\R$.  All four defining properties
are now proved.
\end{proof}

\section{The strict doubled-character model is impossible}

\begin{theorem}[Doubled-Pr\"ufer obstruction]\label{thm:nogo}
Let $G=A\times A$.  Suppose unitaries $\{W_g:g\in G\}$ in a von Neumann
algebra satisfy
\[
 W_gW_h=\sigma(g,h)W_{g+h}\qquad(g,h\in G)
\]
for scalars $\sigma(g,h)\in\T$.  Then all the $W_g$ commute.  In particular,
such a family cannot be an orthonormal basis generating a noncommutative
factor.
\end{theorem}

\begin{proof}
The scalar commutator
\[
 c(g,h)=\frac{\sigma(g,h)}{\sigma(h,g)}
\]
satisfies $W_gW_h=c(g,h)W_hW_g$.  Associativity, or a direct commutator
calculation, shows that $c$ is a bicharacter.  Thus
\[
 \Phi:G\longrightarrow\widehat G,\qquad
 \Phi(g)(h)=c(g,h),
\]
is a group homomorphism.

The group $A=\Z[1/2]/\Z$ is divisible, hence so is $G$.  Its Pontryagin dual
is $\widehat G=\Z_2\times\Z_2$.  The image of a divisible group under a
homomorphism is divisible.  But $\Z_2^2$ has no nonzero divisible subgroup:
if $y$ belongs to such a subgroup, then for every $k$ it is divisible there
by $2^k$, so
\[
 y\in\bigcap_{k\ge0}2^k\Z_2^2=\{0\}.
\]
Therefore $\Phi=0$, so $c(g,h)=1$ for all $g,h$.  The unitaries commute.  If
their linear span were weakly dense, the generated von Neumann algebra would
be abelian, not a factor.
\end{proof}

\begin{corollary}[No nested carry-Weyl bases]\label{cor:finite}
There is no sequence of projective unitary bases of
$M_{2^N}$ indexed by $(\Z/2^N\Z)^2$ which is compatible, up to phases, with
unital matrix-algebra inclusions and with the classical character inclusions
$a\mapsto2a$ in both coordinates.
\end{corollary}

\begin{proof}
A compatible sequence would yield in the inductive limit a projective unitary
family indexed by
\[
 \varinjlim_N(\Z/2^N\Z)^2=A^2.
\]
Each finite projective matrix basis has a nondegenerate commutator pairing:
an element in its radical would give a nonidentity central basis unitary in a
full matrix algebra.  The limit would therefore contain noncommuting
unitaries, contradicting Theorem~\ref{thm:nogo}.

Equivalently, at the finite level every injection into
$(\Z/2^{N+1}\Z)^2$ has image inside (indeed equal to) twice that group.
Pulling an alternating pairing back through two doubled coordinates multiplies
its coefficient by four; relative to denominator $2^{N+1}$ this becomes an
even, hence degenerate, coefficient modulo $2^N$.
\end{proof}

\section{Interpretation and limitations}

Theorem~\ref{thm:positive} gives an explicit affirmative construction with
the basic Fourier-system features used around the source question.  It also
explains why the natural index group is $A\times\Z$: the pairing
$(n,s)\mapsto e^{-2\pi ins}$ between integer odometer translations and
$2$-adic characters is nondegenerate.  Replacing the integer direction by a
second divisible copy of $A$ kills every possible scalar commutator by
Theorem~\ref{thm:nogo}.

The result does \emph{not} assert that some one-dimensional enumeration of
$\{X_{r,n}\}$ has the same Ces\`aro maximal inequalities as the Vilenkin
ordering in \cite{JLZZ}.  Nor can the basis be ordered so that its canonical
finite initial blocks are nested cyclic Weyl bases: Corollary~\ref{cor:finite}
is exactly the obstruction.  Those are analytic or ordering requirements not
stated in Question 5.17.

\section{Verification and novelty status}

The proof uses no computational or unproved dependency.  The multiplication,
trace orthogonality, factor property, and matrix approximation have been
checked independently line by line.  The strict no-go proof reduces to the
elementary fact that $\Z_2^2$ has no nonzero divisible subgroup.

A bounded literature search on 11 August 2026 used the exact question phrase,
the source title and arXiv id, the author names, and combinations of
``noncommutative $2$-adic integer system,'' ``odometer crossed product,''
``Pr\"ufer group twisted group algebra,'' and ``noncommutative solenoid.''  No
paper explicitly answering Question 5.17 was found.  Noncommutative
$p$-solenoids in \cite{LP11,LP13} twist $\Z[1/p]^2$, not the doubled classical
character group $(\Z[1/p]/\Z)^2$, so they are adjacent rather than exact
answers.  The odometer and crossed-product ingredients are standard; the
identification with Question 5.17 and the doubled-Pr\"ufer obstruction were
not found in the bounded search.  Novelty confidence is therefore moderate.

\paragraph{Human-review recommendation.}
Verify (i) that the structural definition captures the authors' intended
undefined term, (ii) the convention in the phase of \eqref{eq:product},
(iii) the nesting identity \eqref{eq:nesting} and $L_2$ approximation of $u$,
and (iv) the claim that every scalar commutator of an $A^2$-projective family
lands in $\widehat{A^2}=\Z_2^2$.  If these checks pass, retain this as a full
affirmative construction plus a full no-go theorem for the strict model.

\begin{thebibliography}{9}
\bibitem{JLZZ}
Y.~Jiao, S.~Luo, T.~Zhao, and D.~Zhou,
\emph{Almost uniform convergence for noncommutative Vilenkin-Fourier series},
arXiv:2506.14431, 2025.

\bibitem{LP11}
F.~Latr\'emoli\`ere and J.~Packer,
\emph{Noncommutative solenoids}, arXiv:1110.6227, 2011.

\bibitem{LP13}
F.~Latr\'emoli\`ere and J.~Packer,
\emph{Noncommutative solenoids and their projective modules},
arXiv:1311.1193, 2013.
\end{thebibliography}

\end{document}
